\chapter{МАТЕМАТИЧЕСКАЯ МОДЕЛЬ И ЧИСЛЕННЫЕ МЕТОДЫ}

В главе формализуется математическая постановка расчёта
термодинамических свойств воды и водяного пара.
При изложении используются базовые термодинамические представления
и инженерные приёмы моделирования,
характерные для учебной литературы по физической газовой динамике~\cite{gidaspov_severina_2014}.
Нормативной основой служит IAPWS-IF97~\cite{iapws_if97_2012}.
Стандарт задаёт региональную топологию, уравнения состояния и часть обратных зависимостей.
Поэтому вычислительный алгоритм включает обязательные шаги
классификации области, выбора модели и контроля применимости.
Принятые здесь соотношения и численные процедуры
непосредственно используются в реализации \texttt{if97\_core},
описанной в следующей главе.

\section{Разбиение области IAPWS-IF97 на регионы}

Рабочая область IF97 разделена на регионы,
в каждом из которых применяются собственные уравнения
и собственные нормировки безразмерных переменных.
Такое разбиение отражает физические различия
между сжатой жидкостью, перегретым паром и околокритической областью.
Одна универсальная аналитическая формула для всей области
в промышленной постановке не обеспечивает одновременно
точность, устойчивость и вычислительную эффективность.

Следовательно, для любого маршрута расчёта
определение региона является обязательным этапом.
Только после классификации состояния корректно
выбирается модель и вычисляются итоговые свойства.

\textbf{Краткая характеристика стандарта IAPWS-IF97.}
IF97 представляет собой промышленную региональную формулировку
для чистой воды и водяного пара,
ориентированную на инженерные расчёты,
где одновременно важны точность,
вычислительная скорость
и устойчивость на границах областей~\cite{iapws_if97_2012}.
Стандарт не задаёт одно универсальное уравнение
	для всей рабочей области,
	а использует набор специализированных аппроксимаций:
	основные уравнения для регионов~1--5,
	аналитическую зависимость для линии насыщения региона~4,
	а также часть обратных зависимостей для входных пар
	$(p, h)$ и $(p, s)$.
При этом IF97 имеет и ограничения:
он относится к чистой воде,
	не предназначен для солевых растворов и смесей,
	не покрывает все возможные обратные постановки
	в аналитической форме,
	а для региона~5 не учитывает эффекты термической диссоциации.
Именно поэтому в вычислительном ядре
аналитические соотношения IF97
дополняются явной валидацией,
маршрутизацией по регионам
и численными процедурами на сложных участках области.

Геометрия рабочей области стандарта наглядно показана на рисунке~\ref{fig:if97_regions}.
Из него видно, что регионы разделены не произвольно, а по физически значимым границам:
линиям насыщения, критической области и высоким температурам,
поэтому выбор формулы в ядре всегда связан с положением точки на этой диаграмме.

\begin{figure}[htbp]
    \centering
    \begin{tikzpicture}
        
        % ЛЕВЫЙ ГРАФИК: Регионы 1, 2, 3 и 4
        \begin{axis}[
            name=ax1,
            width=10cm, height=8cm, % Немного сузили, чтобы точно влезло в поля
            xmin=250, xmax=1150,
            ymin=0, ymax=110,
            axis y line=left,
            axis x line=bottom,
            ylabel={Давление $p$, МПа},
            xlabel={Температура $T$, К},
            xtick={273.15, 623.15, 1073.15},
            ytick={0, 50, 100}, % Оставили только главные метки
            extra y ticks={22.06}, % Добавили критическое давление отдельно
            extra y tick labels={22{,}06},
            grid=both,
            grid style={dashed, gray!30},
            enlargelimits=false,
            ticklabel style={/pgf/number format/use comma},
            clip=false % Отключаем обрезку элементов, выходящих за рамки
        ]
        
        % Регион 4: Линия насыщения (по реперным точкам IF97)
        \addplot [orange, ultra thick, smooth] coordinates {
            (273.15, 0.000611)
            (373.15, 0.1013)
            (473.15, 1.554)
            (573.15, 8.581)
            (647.096, 22.064)
        };
        
        % Граница между Регионами 2 и 3 (Уравнение B23)
        \addplot [domain=623.15:863.15, samples=50, thick, dashed] 
            {13.918837 + 0.001019297 * (x - 572.5446)^2};
            
        % Вертикальная граница между Регионами 1 и 3
        \draw[thick, dashed] (axis cs:623.15, 16.53) -- (axis cs:623.15, 100);
        
        % Вертикальная граница между Регионами 2 и 5
        \draw[thick] (axis cs:1073.15, 0) -- (axis cs:1073.15, 100);
        
        % Верхняя граница графика (p = 100 МПа)
        \draw[thick] (axis cs:273.15, 100) -- (axis cs:1073.15, 100);
        
        % Подписи регионов
        \node at (axis cs: 440, 60) [font=\large] {Регион 1};
        \node at (axis cs: 870, 30) [font=\large] {Регион 2};
        \node at (axis cs: 720, 80) [font=\large] {Регион 3};
        
        % Выноска для Региона 4 (сделали аккуратнее)
        \draw[<-, thick] (axis cs: 573.15, 8.5) -- (axis cs: 750, 15) 
            node[anchor=west, font=\large] {Регион 4};
            
        \end{axis}
        
        % ПРАВЫЙ ГРАФИК: Регион 5 (Разрыв оси)
        \begin{axis}[
            name=ax2,
            at={(ax1.south east)},
            anchor=south west, % Жестко выравниваем левый нижний угол
            xshift=0.8cm, % Увеличили отступ для значка разрыва
            width=3.5cm, height=8cm,
            xmin=2173.15, xmax=2373.15,
            ymin=0, ymax=110,
            axis y line=none, 
            axis x line=bottom,
            xtick={2273.15},
            grid=both,
            grid style={dashed, gray!30},
            enlargelimits=false,
            ticklabel style={/pgf/number format/use comma},
            clip=false % Важно! Чтобы надписи не обрезались правым краем
        ]
        
        % Верхняя и правая границы Региона 5
        \draw[thick] (axis cs:2173.15, 50) -- (axis cs:2273.15, 50);
        \draw[thick] (axis cs:2273.15, 0) -- (axis cs:2273.15, 50);
        
        \node at (axis cs: 2223, 25) [font=\large] {Регион 5};
        \node at (axis cs: 2223, 55) {50 МПа};
        
        \end{axis}
        
        % Отрисовка значка "разрыва" оси X между графиками
        \path (ax1.south east) -- (ax2.south west) coordinate[midway] (mid);
        \draw[thick] (mid) ++(-0.15cm, -0.2cm) -- ++(0.1cm, 0.4cm);
        \draw[thick] (mid) ++(0.05cm, -0.2cm) -- ++(0.1cm, 0.4cm);
        
    \end{tikzpicture}
    \caption{Диаграмма $p$--$T$ с границами регионов IAPWS-IF97}
    \label{fig:if97_regions_fixed}
\end{figure}


\subsection{Характеристика регионов 1--5}

IF97 выделяет пять основных регионов и отдельную область насыщения.
Регионы различаются физическим смыслом
и используемым термодинамическим потенциалом.
Сводная характеристика приведена в таблице~\ref{tab:if97_regions_overview}.

\begin{table}[htbp]
\centering
\caption{Регионы IAPWS-IF97 и используемые формулировки}
\label{tab:if97_regions_overview}
\begin{tblr}{
  width=\linewidth,
  colspec={X[1,l] X[2,l] X[2,l]},
  row{1}={font=\bfseries},
  hline{1,2,Z} = {solid},
}
Регион & Физический смысл & Потенциал и независимые переменные \\
1 & Сжатая вода (подохлаженная жидкость) & Гиббс, $(p, T)$ \\
2 & Перегретый пар & Гиббс, $(p, T)$ \\
3 & Высокоплотная околокритическая область & Гельмгольц, $(\rho, T)$ \\
4 & Линия насыщения и двухфазная область & Насыщение, $(p, x)$ или $(T, x)$ \\
5 & Высокотемпературная область (пар) & Гиббс, $(p, T)$ \\
\end{tblr}
\end{table}

В инженерных сценариях используются входные пары
$(p, T)$, $(p, h)$, $(p, s)$, $(p, x)$ и $(\rho, T)$.
На выходе формируется унифицированное состояние:
$p$, $T$, $v$, $\rho$, $h$, $s$, $u$, $c_p$, $w$ и номер региона IF97.

\paragraph*{Регион 1: сжатая вода.}
Регион~1 описывает стабильную однофазную жидкость
	в диапазоне температур от $273{,}15$ до $623{,}15$~К
	при давлениях от линии насыщения до $100$~МПа~\cite{iapws_if97_2012}.
Для этого региона используется формулировка через потенциал Гиббса,
	а безразмерные переменные нормируются константами
	$p^{*}=16{,}53$~МПа и $T^{*}=1386$~К.
Прямая аппроксимация задаётся суммой из $34$ табличных коэффициентов
$n_i$ и степеней $I_i$, $J_i$.

	Особенностью региона~1 является согласование
с термодинамическим нулём на тройной точке:
коэффициенты уравнения подобраны так,
чтобы удельные внутренняя энергия и энтропия
насыщенной жидкости обращались в нуль.
Стандарт также допускает физически корректное описание
метастабильной жидкости в узкой области,
примыкающей к линии насыщения.

	Для инженерных обратных расчётов
	в регионе~1 предусмотрены аналитические зависимости
		$T(p, h)$ и $T(p, s)$.
Обе зависимости также задаются табличными наборами коэффициентов
по $20$ коэффициентов в каждой аппроксимации.
Это позволяет избегать итераций
в основной жидкофазной части рабочей области.

\paragraph*{Регион 2: перегретый и метастабильный пар.}
Регион~2 описывает перегретый пар
	и газообразное состояние воды
	в трёх температурно-барических диапазонах:
	при $273{,}15 \le T \le 623{,}15$~К и $0 < p \le p_{\mathrm{sat}}(T)$,
	при $623{,}15 < T \le 863{,}15$~К и $0 < p \le p_{B23}(T)$,
	а также при $863{,}15 < T \le 1073{,}15$~К и $0 < p \le 100$~МПа
	~\cite{iapws_if97_2012}.

	Математически регион~2 задаётся через энергию Гиббса,
	разбитую на идеально-газовую и остаточную части.
	В качестве нормировочных констант
	используются $p^{*}=1$~МПа и $T^{*}=540$~К.
Идеально-газовая часть аппроксимации
содержит $9$ табличных коэффициентов,
остаточная часть --- $43$ коэффициента.
	Отдельно стандарт учитывает метастабильный пар
	в примыкающей к линии насыщения области;
	для диапазона давлений до $10$~МПа
	используется специальная модификация остаточной части.

Для обратных расчётов
регион дополнительно разбивается на подобласти 2a, 2b и 2c.
Это разбиение используется только для аналитических зависимостей
$T(p, h)$ и $T(p, s)$,
которые задаются отдельными табличными наборами коэффициентов
для каждой подобласти,
и ускоряют расчёты в паровой области
без внешнего итерационного контура.

\paragraph*{Регион 3: высокоплотная околокритическая область.}
Регион~3 охватывает область
	от $623{,}15$ до $863{,}15$~К
	между границей $B23$ и верхним пределом давления $100$~МПа
	~\cite{iapws_if97_2012}.
	Именно здесь расположены критическая точка воды
	с параметрами $T_{c}=647{,}096$~К,
	$p_{c}=22{,}064$~МПа
	и $\rho_{c}=322$~кг/м\textsuperscript{3},
	а также наиболее выраженные нелинейности термодинамических свойств.

	В отличие от регионов~1, 2 и~5,
	для региона~3 независимыми переменными служат плотность и температура.
Поэтому стандарт использует не потенциал Гиббса,
а безразмерную энергию Гельмгольца,
нормированную критическими значениями
$\rho^{*}=\rho_{c}$ и $T^{*}=T_{c}$.
Такой выбор естественен для околокритической области,
где давление удобнее рассматривать
как производную по плотности,
а не как независимый аргумент.
	Базовая аппроксимация региона~3
	использует $40$ табличных коэффициентов
	$n_i$, $I_i$, $J_i$.

	Для региона~3 в IF97 не вводятся аналитические обратные уравнения
		вида $T(p, h)$ или $T(p, s)$.
Поэтому при постановках $(p, T)$, $(p, h)$ и $(p, s)$
требуется либо использование специальных вспомогательных релизов,
либо численное решение нелинейных уравнений.
Именно эта особенность определяет необходимость
дополнительных численных методов в проекте.

\paragraph*{Регион 4: линия насыщения и двухфазная область.}
Регион~4 описывает не двумерную область состояния,
	а кривую насыщения от тройной точки до критической точки:
	$273{,}15 \le T \le 647{,}096$~К
	и $611{,}213$~Па $\le p \le 22{,}064$~МПа
	~\cite{iapws_if97_2012}.
Для него не вводится собственное фундаментальное уравнение
в виде энергии Гиббса или Гельмгольца.
Вместо этого используются аналитически согласованные зависимости
$p_{\mathrm{sat}}(T)$ и $T_{\mathrm{sat}}(p)$,
получаемые из одного неявного полиномиального соотношения.
Это соотношение содержит $10$ табличных констант
$n_1,\dots,n_{10}$.

	Полные свойства на линии насыщения
	вычисляются через соседние однофазные регионы.
	Ниже $623{,}15$~К насыщенная жидкость
	и насыщенный пар вычисляются по формулам регионов~1 и~2,
	а в верхней околокритической части линии насыщения
	используется формулировка региона~3.
Внутри двухфазной области состояние дополнительно задаётся
степенью сухости $x$,
по которой выполняется расчёт свойств смеси.

\paragraph*{Регион 5: высокотемпературная область.}
Регион~5 используется для высокотемпературного пара
	при $1073{,}15 < T \le 2273{,}15$~К
	и давлениях до $50$~МПа~\cite{iapws_if97_2012}.
	Как и регион~2,
	он описывается через сумму идеально-газовой
	и остаточной частей энергии Гиббса.
	Для нормировки применяются
	$p^{*}=1$~МПа и $T^{*}=1000$~К.
Идеально-газовая и остаточная части
содержат по $6$ табличных коэффициентов.

	Для региона~5 в стандарте не предусмотрены
	аналитические обратные уравнения.
Поэтому при расчётах по энтальпии или энтропии
температура в этой области
должна восстанавливаться численными методами.
Кроме того, модель относится к химически чистой,
недиссоциированной воде
и не учитывает эффекты термической диссоциации,
которые при столь высоких температурах
могут становиться значимыми в специальных приложениях.

\subsection{Уравнения линии насыщения и границы B23}

	Линия насыщения (регион~4) является одномерным множеством состояний,
в которых жидкая и паровая фазы находятся в равновесии.
Физически это означает,
что по обе стороны линии располагаются однофазные области,
а внутри купола насыщения существует двухфазная смесь.
Поэтому для входной пары $(p, T)$
оба параметра на линии насыщения уже не являются независимыми:
задание одного из них однозначно определяет второй.
Формально линия насыщения соответствует условию фазового равновесия,
при котором удельные энергии Гиббса жидкой и паровой фаз совпадают.

\textbf{Тройная и критическая точки.}
	Нижней опорной точкой термодинамической шкалы
	служит тройная точка воды
	с параметрами $T_{t}=273{,}16$~К и $p_{t}=611{,}657$~Па,
	в которой сосуществуют твёрдая, жидкая и паровая фазы.
	Именно к этой точке привязана нормировка региона~1,
где удельные внутренняя энергия и энтропия
насыщенной жидкости принимаются равными нулю.
	Верхней точкой линии насыщения
	служит критическая точка
	$T_{c}=647{,}096$~К,
	$p_{c}=22{,}064$~МПа,
	$\rho_{c}=322$~кг/м\textsuperscript{3}.
В критической точке различие между насыщенной жидкостью
и насыщенным паром исчезает,
а разность их удельных объёмов стремится к нулю.
При этом выполняются условия
\begin{equation}
\left(\frac{\partial p}{\partial \rho}\right)_{T} = 0,
\qquad
\left(\frac{\partial^{2} p}{\partial \rho^{2}}\right)_{T} = 0,
\end{equation}
что и делает околокритическую область
особенно чувствительной к численным ошибкам.

\textbf{Аналитическое задание линии насыщения.}
Для региона~4 стандарт IF97 использует единое неявное полиномиальное соотношение
~\cite{iapws_if97_2012}:
\begin{equation}
\beta^{2}\vartheta^{2}
+ n_{1}\beta^{2}\vartheta
+ n_{2}\beta^{2}
+ n_{3}\beta\vartheta^{2}
+ n_{4}\beta\vartheta
+ n_{5}\beta
+ n_{6}\vartheta^{2}
+ n_{7}\vartheta
+ n_{8}
= 0,
\end{equation}
\begin{gostwhere}
  \item $\beta$~--- безразмерная переменная давления, $\beta = \left(\frac{p_{\mathrm{sat}}}{p^{*}}\right)^{1/4}$;
  \item $\vartheta$~--- вспомогательная безразмерная переменная температуры, $\vartheta = \frac{T}{T^{*}} + \frac{n_{9}}{\frac{T}{T^{*}} - n_{10}}$;
  \item $p^{*}$~--- нормировочное давление, $p^{*}=1~\mbox{МПа}$;
  \item $T^{*}$~--- нормировочная температура, $T^{*}=1~\mbox{К}$.
\end{gostwhere}

Все коэффициенты $n_{1},\dots,n_{10}$
задаются стандартом в табличной форме.

Из этого полинома выводится прямое уравнение насыщения
для зависимости $p_{\mathrm{sat}}(T)$:
\begin{equation}
p_{\mathrm{sat}}(T) = p^{*}
\left[
\frac{2C}{-B + \sqrt{B^{2} - 4AC}}
\right]^{4},
\end{equation}
\begin{gostwhere}
  \item $A$~--- вспомогательный коэффициент, $A = \vartheta^{2} + n_{1}\vartheta + n_{2}$;
  \item $B$~--- вспомогательный коэффициент, $B = n_{3}\vartheta^{2} + n_{4}\vartheta + n_{5}$;
  \item $C$~--- вспомогательный коэффициент, $C = n_{6}\vartheta^{2} + n_{7}\vartheta + n_{8}$.
\end{gostwhere}

Аналогично получается и обратное уравнение
для зависимости $T_{\mathrm{sat}}(p)$:
\begin{equation}
T_{\mathrm{sat}}(p) = \frac{T^{*}}{2}
\left[
n_{10} + D - \sqrt{(n_{10}+D)^{2} - 4(n_{9} + n_{10}D)}
\right],
\end{equation}
\begin{gostwhere}
  \item $D$~--- вспомогательная величина, $D = \frac{2G}{-F - \sqrt{F^{2} - 4EG}}$;
  \item $E$~--- вспомогательный коэффициент, $E = \beta^{2} + n_{3}\beta + n_{6}$;
  \item $F$~--- вспомогательный коэффициент, $F = n_{1}\beta^{2} + n_{4}\beta + n_{7}$;
  \item $G$~--- вспомогательный коэффициент, $G = n_{2}\beta^{2} + n_{5}\beta + n_{8}$.
\end{gostwhere}

Таким образом, линия насыщения
задаётся аналитически,
а переход между постановками $(T)$ и $(p)$
не требует итерационного решения.

\textbf{Почему регион~4 является численно сложной областью.}
Внутри двухфазной области свойства среды
не задаются одним фундаментальным уравнением состояния.
Вместо этого требуется вычислить свойства насыщенной жидкости
и насыщенного пара,
а затем выполнить смешение по степени сухости $x$.
Кроме того, на линии насыщения и вблизи критической точки
производные свойств ведут себя жёстко:
изобарная теплоёмкость быстро возрастает,
скорость звука демонстрирует выраженный провал,
а величина $v''-v'$ убывает,
что ухудшает обусловленность обратных задач.

\textbf{Как эта сложность обрабатывается в модели.}
В проекте проблемы региона~4 решаются комбинированно.
Сама линия насыщения вычисляется по аналитическим формулам
$p_{\mathrm{sat}}(T)$ и $T_{\mathrm{sat}}(p)$.
Свойства насыщенной жидкости и насыщенного пара
ниже $623{,}15$~К рассчитываются по формулам регионов~1 и~2,
а в околокритической части линии насыщения ---
через формулировку региона~3.
Для смеси используются явные правила смешения по $x$,
а при вычислениях на границах вводятся проверки диапазонов,
контроль принадлежности куполу насыщения
и отдельная обработка плохо обусловленного поведения
вблизи критической точки.

Граница B23 отделяет регион~2 от региона~3.
Она определяет момент перехода
от формулировки через потенциал Гиббса
к формулировке через потенциал Гельмгольца.
В безразмерной записи эта граница задаётся
квадратичной аппроксимацией вида
\begin{equation}
\pi = n_{1} + n_{2}\theta + n_{3}\theta^{2},
\end{equation}
\begin{gostwhere}
  \item $\pi$~--- безразмерное давление, $\pi = p/p^{*}$;
  \item $\theta$~--- безразмерная температура, $\theta = T/T^{*}$~\cite{iapws_if97_2012}.
\end{gostwhere}

Корректная обработка B23 обязательна
для маршрутизации в постановках $(p, T)$, $(p, h)$ и $(p, s)$.

\section{Термодинамические потенциалы и расчётные зависимости}

Модель IF97 вычисляет свойства через производные
безразмерных термодинамических потенциалов.
Такое представление обеспечивает единую форму записи,
корректное масштабирование переменных
и лучшую численную устойчивость.

В регионах 1, 2 и 5 используется потенциал Гиббса.
В регионе 3 используется потенциал Гельмгольца.
Регион~4 описывается зависимостями насыщения и правилами смешения.

\subsection{Свободная энергия Гиббса для регионов 1, 2 и 5}

Для регионов 1, 2 и 5 стандарт IF97 задаёт безразмерную энергию Гиббса
\begin{equation}
\gamma(\pi, \tau) = \frac{g(p, T)}{R T},
\end{equation}
\begin{gostwhere}
  \item $g$~--- удельная энергия Гиббса;
  \item $R$~--- газовая постоянная воды и водяного пара.
\end{gostwhere}

Безразмерные переменные определяются как
\begin{equation}
\pi = \frac{p}{p^{*}}, \qquad \tau = \frac{T^{*}}{T}.
\end{equation}
Нормировочные константы $p^{*}$ и $T^{*}$ задаются стандартом отдельно для каждого региона.

Для региона~1 функция $\gamma$ задаётся единым полиномиальным разложением
со смещёнными аргументами:
\begin{equation}
\gamma(\pi,\tau) = \sum_{i} n_{i}\,(7{,}1-\pi)^{I_{i}}\,(\tau-1{,}222)^{J_{i}}.
\end{equation}
В этой области используются нормировки
$p^{*}=16{,}53$~МПа и $T^{*}=1386$~К.
Здесь аппроксимация использует $34$ табличных коэффициента.

Для регионов~2 и~5 применяется представление
\begin{equation}
\gamma(\pi,\tau) = \gamma^{o}(\pi,\tau) + \gamma^{r}(\pi,\tau),
\end{equation}
в котором $\gamma^{o}$ описывает идеально-газовую часть,
а $\gamma^{r}$ --- поправку на реальное взаимодействие молекул.
Для региона~2 используются
$p^{*}=1$~МПа и $T^{*}=540$~К,
а для региона~5 ---
$p^{*}=1$~МПа и $T^{*}=1000$~К.
Для региона~2 табличные массивы содержат
$9$ коэффициентов в идеально-газовой части
и $43$ коэффициента в остаточной части.
Для региона~5 используются по $6$ коэффициентов
в идеально-газовой и остаточной частях.

Свойства выражаются через частные производные $\gamma(\pi,\tau)$.
Используются обозначения
$\gamma_{\pi} = \partial \gamma/\partial \pi$,
$\gamma_{\tau} = \partial \gamma/\partial \tau$,
$\gamma_{\pi\pi} = \partial^{2}\gamma/\partial \pi^{2}$,
$\gamma_{\tau\tau} = \partial^{2}\gamma/\partial \tau^{2}$,
$\gamma_{\pi\tau} = \partial^{2}\gamma/\partial \pi \partial \tau$.

Ключевые зависимости, используемые в вычислениях, включают:
\begin{gather}
v = \frac{R T}{p}\,\pi\,\gamma_{\pi}, \\
\rho = \frac{1}{v} = \frac{p}{R T\,\pi\,\gamma_{\pi}}, \\
h = R T\,\tau\,\gamma_{\tau}, \\
s = R \left(\tau\,\gamma_{\tau} - \gamma\right), \\
u = R T \left(\tau\,\gamma_{\tau} - \pi\,\gamma_{\pi}\right), \\
c_{p} = -R\,\tau^{2}\,\gamma_{\tau\tau}, \\
w = \sqrt{R T\,
\frac{\gamma_{\pi}^{2}}
{\frac{\left(\gamma_{\pi} - \tau\gamma_{\pi\tau}\right)^{2}}
{\tau^{2}\gamma_{\tau\tau}}
- \gamma_{\pi\pi}} }.
\end{gather}
Именно эти формулы образуют базовый набор
для прямых расчётов по входной паре $(p, T)$
в регионах~1, 2 и~5~\cite{iapws_if97_2012}.

\subsection{Свободная энергия Гельмгольца для региона 3}

Для региона~3 стандарт использует формулировку через безразмерную энергию Гельмгольца
\begin{equation}
\phi(\delta, \tau) = \frac{f(\rho, T)}{R T},
\end{equation}
\begin{gostwhere}
  \item $f$~--- удельная энергия Гельмгольца.
\end{gostwhere}

Безразмерные переменные определяются как
\begin{equation}
\delta = \frac{\rho}{\rho^{*}}, \qquad \tau = \frac{T^{*}}{T}.
\end{equation}
Для региона~3 фундаментальная зависимость имеет вид
\begin{equation}
\phi(\delta,\tau) = n_{1}\ln\delta + \sum_{i=2}^{40} n_{i}\,\delta^{I_{i}}\,\tau^{J_{i}},
\end{equation}
а нормировочные константы совпадают
с критическими параметрами воды:
$\rho^{*}=\rho_{c}=322$~кг/м\textsuperscript{3} и $T^{*}=T_{c}=647{,}096$~К.

В этой формулировке давление выражается через производную $\phi$ по плотности:
\begin{equation}
p = \rho R T\,\delta\,\phi_{\delta}.
\end{equation}
Это соотношение показывает,
что для региона~3 естественными независимыми переменными
являются $\rho$ и $T$.
При постановке $(p, T)$ требуется сначала определить плотность,
согласованную с заданным давлением.

После определения $\rho$
остальные свойства вычисляются через производные $\phi$:
\begin{gather}
v = \frac{1}{\rho} = \frac{1}{\delta\,\rho^{*}}, \\
u = R T\,\tau\,\phi_{\tau}, \\
s = R \left(\tau\,\phi_{\tau} - \phi\right), \\
h = R T \left(\tau\,\phi_{\tau} + \delta\,\phi_{\delta}\right), \\
c_{p} = R \left[
-\tau^{2}\phi_{\tau\tau}
+ \frac{\left(\delta\phi_{\delta} - \delta\tau\phi_{\delta\tau}\right)^{2}}
{2\delta\phi_{\delta} + \delta^{2}\phi_{\delta\delta}}
\right], \\
w = \sqrt{R T \left[
2\delta\phi_{\delta} + \delta^{2}\phi_{\delta\delta}
- \frac{\left(\delta\phi_{\delta} - \delta\tau\phi_{\delta\tau}\right)^{2}}
{\tau^{2}\phi_{\tau\tau}}
\right]}.
\end{gather}
Формально это такие же табличные аппроксимации,
но в пространстве переменных $(\rho, T)$,
что и делает регион~3 естественным для расчётов
по высокой плотности и околокритическим состояниям~\cite{iapws_if97_2012}.
Важной особенностью является отсутствие
в стандарте аналитических обратных зависимостей для региона~3,
поэтому обратные постановки в этой области
опираются на численные решатели.

\section{Прямые и обратные маршруты расчёта}

Практическая применимость модели определяется набором маршрутов расчёта.
Маршрут задаётся входной парой параметров
и определяет последовательность вычислительных шагов.
В проекте используются как прямые зависимости IF97,
так и обратные постановки с аналитическими и численными процедурами.

\subsection{Прямой маршрут по параметрам \texorpdfstring{$(p, T)$ и $(\rho, T)$}{(p, T) и (rho, T)}}

Прямой маршрут $(p, T)$ является базовым для IF97.
Для него, а также для естественной для региона~3 постановки $(\rho, T)$,
расчёт строится вокруг прямого использования фундаментальных уравнений состояния.
По входной паре $(p, T)$ выполняются следующие этапы:
\begin{enumerate}
\item валидация диапазонов входных значений в пределах области применимости стандарта;
\item определение региона IF97 по топологическим границам (линия насыщения, граница B23 и др.);
\item вычисление безразмерных переменных и значений потенциала в выбранном регионе;
\item расчёт производных потенциала и итоговых свойств;
\item формирование унифицированного результата, включающего свойства и номер региона.
\end{enumerate}

В регионе~3 маршрут $(p, T)$ включает шаг определения плотности,
поскольку уравнения заданы в переменных $(\rho, T)$.
Для этого используются обратные зависимости $v(p, T)$,
опубликованные IAPWS в дополнительном релизе~\cite{iapws_vpt_region3},
и контроль физической допустимости найденного состояния.

Если же входная пара сразу задана как $(\rho, T)$,
то для региона~3 свойства вычисляются непосредственно
через потенциал Гельмгольца без дополнительной смены переменных.
В однофазных регионах 1, 2 и 5
маршрут $(\rho, T)$ трактуется как расширенная инженерная постановка:
по заданным плотности и температуре сначала восстанавливается давление,
после чего вычисление продолжается по стандартному маршруту $(p, T)$.

\subsection{Итерационные методы для маршрутов \texorpdfstring{$(p, h)$, $(p, s)$ и $(p, x)$}{(p, h), (p, s) и (p, x)}}

Обратные маршруты применяются,
когда состояние задаётся энергетическими параметрами.
К ним относятся входные пары $(p, h)$ и $(p, s)$.
В этих постановках нужно восстановить температуру
и далее вычислить полный набор свойств.

Для регионов и подрегионов,
где IF97 предоставляет аналитические обратные формулы,
температура определяется без итераций.
Это снижает вычислительные затраты
в массовых инженерных сценариях.

В регионе~1 используются единые зависимости
$T(p, h)$ и $T(p, s)$.
В регионе~2 стандарт задаёт отдельные семейства формул
для подобластей 2a, 2b и 2c.
Это позволяет учитывать сложную форму паровой области
без полного итерационного решения.
Для региона~5 аналитические обратные уравнения
в IF97 отсутствуют,
поэтому при попадании состояния в высокотемпературную область
требуется численное восстановление неизвестной температуры.

Если постановка $(p, h)$ или $(p, s)$ попадает в регион~3,
возникает система нелинейных уравнений относительно $(\rho, T)$.
В проекте для таких 2D-задач применяется решатель
Левенберга--Марквардта с предварительным сеточным поиском
и демпфированием шага.
Этот подход выбран для устойчивой сходимости
в околокритической области.

Маршрут $(p, x)$ используется для двухфазной области.
Хотя для него не требуется внешний итерационный контур,
он рассматривается вместе с обратными маршрутами,
поскольку также восстанавливает полное состояние
по неполному набору независимых параметров.
По заданным давлению $p$ и степени сухости $x$ определяется температура насыщения $T_{\mathrm{sat}}(p)$ и свойства насыщенной жидкости
($x=0$) и насыщенного пара ($x=1$).
В нижней части линии насыщения
эти свойства берутся из регионов~1 и~2,
а в околокритической части ---
из формулировки региона~3.
Для величин, аддитивных по массе, применяются правила смешения:
\begin{equation}
y = (1-x)\,y_{f} + x\,y_{g},
\end{equation}
\begin{gostwhere}
\item индекс $f$ соответствует насыщенной жидкости;
\item индекс $g$ --- насыщенному пару.
\end{gostwhere}

Удельный объём смеси вычисляется по этой формуле, а плотность определяется как $\rho = 1/v$.
Свойства двухфазной области, включая поведение $c_p$ и $w$,
требуют отдельной интерпретации в прикладных расчётах.

\section{Численные методы расширения модели}

IF97 задаёт основной набор формул и часть обратных зависимостей,
но для полной рабочей области требуются дополнительные численные процедуры.
В проекте они используются как расширение,
сохраняющее физическую основу стандарта.

\subsection[Метод секущих для региона~3 и однофазных областей]{Метод секущих для восстановления плотности в регионе~3 и давления в однофазных областях}

В регионе~3 естественной входной парой является $(\rho, T)$.
Следовательно, при расчёте по входной паре $(p, T)$
сначала требуется восстановить плотность,
согласованную с заданными давлением и температурой.
Эта задача сводится к поиску корня уравнения
\begin{equation}
F_{\rho}(\rho) = p_{\mathrm{if97}}(\rho, T_{\mathrm{target}}) - p_{\mathrm{target}} = 0.
\end{equation}

Общий ход одномерного итерационного контура показан на рисунке~\ref{fig:secant_algorithm}.
Схема отражает последовательность выбора начального приближения,
вычисления невязки, проверки критерия остановки
и возврата к следующей итерации с учётом физических ограничений области.

\begin{figure}[htbp]
    \centering
    \begin{tikzpicture}[
        scale=0.75,
        transform shape,
        node distance=0.8cm and 1.5cm,
        auto,
        block/.style={rectangle, draw, fill=blue!10, text width=5.7cm, align=center, rounded corners, minimum height=0.8cm},
        decision/.style={diamond, draw, fill=orange!12, text width=3.2cm, align=center, inner sep=1pt, aspect=2.1},
        line/.style={draw, -Latex}
    ]
        \node[block] (start) {Входные данные: $p_{\mathrm{target}}$, $T_{\mathrm{target}}$\\или $\rho_{\mathrm{target}}$, $T_{\mathrm{target}}$};
        \node[block, below=of start] (init) {Выбор $p_{0}$ или $\rho_{0}$\\ и формирование второго приближения\\ $p_{1}=p_{0}\pm 1\%$};
        \node[block, below=of init] (calc) {Вызов вычислительного шага:\\
        $f_{0}=F(p_{0})$ или $F(\rho_{0})$\\
        $f_{1}=F(p_{1})$ или $F(\rho_{1})$};
        \node[decision, below=of calc] (tol) {$|f_{1}| \le \varepsilon$\\или\\$|f_{1}-f_{0}| \approx 0$?};
        
        \node[decision, below=of tol] (limit) {$iter \ge iter_{\max}$?};
        
        \node[block, below=of limit] (step) {Расчёт шага метода секущих:\\
        $\Delta x = -f_{1}\dfrac{x_{1}-x_{0}}{f_{1}-f_{0}}$};
        \node[block, below=of step] (clamp) {Ограничение шага и проверка\\ фазовых границ, границ региона\\ и допустимого диапазона};
        \node[block, below=of clamp] (update) {Обновление приближения:\\
        $x_{0}\leftarrow x_{1}$\\
        $x_{1}\leftarrow x_{1}+\Delta x$\\
        $iter \leftarrow iter + 1$};
        
        \node[block, fill=green!18, right=of tol] (ok) {Успешный выход:\\
        возврат согласованного состояния};
        \node[block, fill=red!15, right=of limit] (fail) {Ошибка:\\
        превышен лимит итераций\\или нарушена сходимость};

        \path[line] (start) -- (init);
        \path[line] (init) -- (calc);
        \path[line] (calc) -- (tol);
        \path[line] (tol) -- node[near start, above] {Да} (ok);
        \path[line] (tol) -- node[near start, right] {Нет} (limit);
        \path[line] (limit) -- node[near start, above] {Да} (fail);
        \path[line] (limit) -- node[near start, right] {Нет} (step);
        \path[line] (step) -- (clamp);
        \path[line] (clamp) -- (update);
        \draw[line] (update.west) -- ++(-1.4,0) |- (calc.west);
    \end{tikzpicture}
    \caption{Блок-схема одномерного итерационного контура метода секущих с контролем физических ограничений}
    \label{fig:secant_algorithm}
\end{figure}

Для одномерного восстановления используется метод секущих:
\begin{equation}
\rho_{k+1} = \rho_{k}
- F_{\rho}(\rho_{k})\,
\frac{\rho_{k}-\rho_{k-1}}
{F_{\rho}(\rho_{k})-F_{\rho}(\rho_{k-1})}.
\end{equation}
Начальные приближения $\rho_{0}$ и $\rho_{1}$
выбираются в физически допустимом интервале региона~3,
а итерации завершаются,
когда выполняется условие
\begin{equation}
|F_{\rho}(\rho_{k})| < \varepsilon_{p}
\quad \mbox{или} \quad
|\rho_{k}-\rho_{k-1}| < \varepsilon_{\rho}.
\end{equation}
После восстановления плотности
состояние рассчитывается по прямым формулам региона~3.
Опубликованные IAPWS зависимости $v(p, T)$ для региона~3
могут использоваться как опорная аппроксимация
или как источник для проверки результата~\cite{iapws_vpt_region3}.

Для однофазных регионов~1, 2 и~5
естественной входной парой является $(p, T)$.
Поэтому маршрут $(\rho, T)$
вне двухфазной области удобно трактовать симметрично:
при заданных $\rho_{\mathrm{target}}$ и $T_{\mathrm{target}}$
сначала восстанавливается давление,
после чего расчёт продолжается стандартным маршрутом $(p, T)$.
Целевая функция имеет вид
\begin{equation}
F_{p}(p) = \rho_{\mathrm{if97}}(p, T_{\mathrm{target}}) - \rho_{\mathrm{target}}.
\end{equation}

Итерационный шаг метода секущих записывается как
\begin{equation}
p_{k+1} = p_{k}
- F_{p}(p_{k})\,
\frac{p_{k}-p_{k-1}}
{F_{p}(p_{k})-F_{p}(p_{k-1})}.
\end{equation}
Критерии остановки выбираются по той же схеме,
что и для восстановления плотности:
\begin{equation}
|F_{p}(p_{k})| < \varepsilon_{\rho}
\quad \mbox{или} \quad
|p_{k}-p_{k-1}| < \varepsilon_{p}.
\end{equation}
После нахождения давления
срабатывает обычная маршрутизация IF97,
и вычисление продолжается как расчёт по входной паре $(p, T)$.
Тем самым маршрут $(\rho, T)$
остаётся встроенным в общую архитектуру модели
и не требует отдельного набора формул
для каждого однофазного региона.

\textbf{Многомерные обратные задачи для пар $(p, h)$ и $(p, s)$ в регионе~3.}
Если состояние в регионе~3 задаётся входной парой $(p, h)$
или $(p, s)$,
задача уже не сводится к одномерному поиску.
Неизвестными становятся сразу две переменные:
\begin{equation}
\mathbf{z} =
\left[
\begin{array}{c}
\rho \\
T
\end{array}
\right],
\qquad
\Phi(\mathbf{z}) = \frac{1}{2}\|\mathbf{r}(\mathbf{z})\|_{2}^{2}.
\end{equation}
Для пары $(p, h)$ вектор невязки имеет вид
\begin{equation}
\mathbf{r}(\mathbf{z}) =
\left[
\begin{array}{c}
p_{\mathrm{if97}}(\rho, T) - p_{\mathrm{target}} \\
h_{\mathrm{if97}}(\rho, T) - h_{\mathrm{target}}
\end{array}
\right],
\end{equation}
а для пары $(p, s)$
во второй компоненте используется
$s_{\mathrm{if97}}(\rho, T) - s_{\mathrm{target}}$.

Метод Ньютона--Рафсона
линеаризует систему в окрестности текущего приближения
и на каждом шаге решает задачу
\begin{equation}
\begin{aligned}
J_{k}\,\Delta \mathbf{z}_{k} &= -\mathbf{r}_{k}, \\
\mathbf{z}_{k+1} &= \mathbf{z}_{k} + \Delta \mathbf{z}_{k},
\end{aligned}
\end{equation}
\begin{gostwhere}
\item $J_{k} = \partial \mathbf{r}/\partial \mathbf{z}$ --- матрица Якоби.
\end{gostwhere}

При хорошем начальном приближении
метод обеспечивает быструю локальную сходимость.
Однако в околокритической области
он оказался недостаточно устойчивым:
из-за роста сжимаемости,
уплощения изотерм и вырождения производных давления по плотности
матрица $J_{k}$ становится плохо обусловленной,
а полный шаг $\Delta \mathbf{z}_{k}$
может выбрасывать траекторию за границы допустимого региона
или переводить её на неверную фазовую ветвь.

Поэтому применяется
схема Левенберга--Марквардта:
внутренний цикл решателя строится по тому же принципу,
но работает уже с вектором невязки, матрицей Якоби
и адаптивным демпфированием шага.
\begin{equation}
\left(J_{k}^{\mathsf{T}}J_{k} + \lambda_{k} I\right)\Delta \mathbf{z}_{k}
= -J_{k}^{\mathsf{T}}\mathbf{r}_{k},
\end{equation}
\begin{equation}
\mathbf{z}_{k+1} = \mathbf{z}_{k} + \alpha_{k}\Delta \mathbf{z}_{k},
\qquad
0 < \alpha_{k} \le 1.
\end{equation}
Здесь параметр $\lambda_{k}$ выполняет роль регуляризации:
при больших значениях метод ведёт себя ближе к градиентному спуску,
а при малых --- приближается к гаусс-ньютоновскому шагу.
Демпфирующий множитель $\alpha_{k}$
ограничивает длину шага
и предотвращает перескок через фазовые границы.

В реализованной модификации используются три дополнительных приёма.
Во-первых, начальное приближение выбирается
не произвольно,
а по результатам предварительного сеточного поиска:
\begin{equation}
\mathbf{z}_{0} = \arg\min_{\mathbf{z}\in\mathcal{G}}
\|\mathbf{r}(\mathbf{z})\|_{2},
\end{equation}
\begin{gostwhere}
\item $\mathcal{G}$ --- грубая сетка в допустимой области региона~3.
\end{gostwhere}

Во-вторых, параметр $\lambda_{k}$ адаптируется по качеству шага:
если норма невязки уменьшается,
то $\lambda_{k}$ снижается,
а если шаг оказывается неудачным ---
увеличивается.
В-третьих, для подбора $\alpha_{k}$
используется демпфирование шага,
например по правилу обратного деления:
\begin{equation}
\begin{aligned}
\alpha_{k} &= 2^{-m_{k}}, \\
m_{k} &= \min\left\{
m \ge 0 :
\|\mathbf{r}(\mathbf{z}_{k} + 2^{-m}\Delta \mathbf{z}_{k})\|_{2}
\right. \\
&\qquad \left.
<
\|\mathbf{r}(\mathbf{z}_{k})\|_{2}
\right\}.
\end{aligned}
\end{equation}
Совместно эти модификации решают три разные проблемы:
сеточный поиск уменьшает зависимость от стартовой точки,
регуляризация стабилизирует почти вырожденный Якобиан,
а демпфирование предотвращает осцилляции и выход за пределы региона~3.

На рисунке~\ref{fig:convergence_comparison} показана характерная динамика
сходимости невязки для двух типовых режимов работы решателя.
В однофазной области, удалённой от линии насыщения,
метод Ньютона--Рафсона достигает целевой точности
за несколько итераций,
тогда как в околокритической области
тот же контур переходит в демпфированный режим
и требует заметно большего числа шагов.
Именно это поведение определяет практическую структуру
итеративных решателей в проекте:
в стабильных областях используется быстрый локальный шаг,
а в сложных зонах включаются line search,
регуляризация и контроль региона.

\begin{figure}[htbp]
    \centering
    \begin{tikzpicture}
    \begin{axis}[
        width=\textwidth,
        height=8cm,
        ymode=log,
        ymin=1e-9, ymax=1e1,
        xmin=0, xmax=15,
        xlabel={Номер итерации, $N$},
        ylabel={Абсолютная невязка, $|f(x)|$},
        grid=major,
        grid style={dashed, gray!50},
        legend pos=north east,
        legend style={cells={anchor=west}, font=\footnotesize},
        mark options={solid}
    ]
        \addplot [domain=0:15, dashed, thick, red] {1e-7};
        \addlegendentry{Целевая точность ($10^{-7}$)}

        \addplot[
            color=blue,
            mark=square*,
            thick
        ] coordinates {
            (1, 2.5e0)
            (2, 4.1e-1)
            (3, 1.2e-2)
            (4, 3.5e-5)
            (5, 1.8e-9)
        };
        \addlegendentry{Регион 5 (однофазная область, 1D)}

        \addplot[
            color=green!60!black,
            mark=*,
            thick
        ] coordinates {
            (1, 3.0e0)
            (2, 1.5e0)
            (3, 8.0e-1)
            (4, 4.5e-1)
            (5, 2.0e-1)
            (6, 8.5e-2)
            (7, 3.1e-2)
            (8, 9.0e-3)
            (9, 1.5e-3)
            (10, 2.2e-4)
            (11, 1.1e-5)
            (12, 4.0e-7)
            (13, 1.2e-8)
        };
        \addlegendentry{Регион 3 (околокритическая область, 2D)}

    \end{axis}
    \end{tikzpicture}
    \caption{Сравнительная динамика сходимости итерационных решателей в однофазной и околокритической областях IF97}
    \label{fig:convergence_comparison}
\end{figure}

\subsection{Особенности сходимости вблизи критической точки}

Наиболее сложные численные режимы
возникают вблизи критической точки
и на примыкающей к ней двухфазной границе.
В двухфазной области (регион~4)
температура однозначно задаёт давление насыщения:
\begin{equation}
p = p_{\mathrm{sat}}(T).
\end{equation}
Следовательно, при задании входной пары $(\rho, T)$
дополнительной неизвестной становится не давление,
а степень сухости $x$.

Связь между заданной плотностью смеси
и степенью сухости удобно записывать
через удельный объём.
Для насыщенной жидкости и насыщенного пара
обозначим свойства через штрихи:
$v'$ и $v''$, либо эквивалентно $\rho'$ и $\rho''$.
Тогда для двухфазной смеси выполняется соотношение
\begin{equation}
\begin{aligned}
v &= (1-x)\,v' + x\,v'', \\
v &= \frac{1}{\rho_{\mathrm{target}}}.
\end{aligned}
\end{equation}
Эквивалентная запись через плотности имеет вид
\begin{equation}
\frac{1}{\rho_{\mathrm{target}}}
=
\frac{1-x}{\rho'}
+
\frac{x}{\rho''}.
\end{equation}
Отсюда степень сухости определяется явно:
\begin{equation}
\begin{aligned}
x &= \frac{v - v'}{v'' - v'} \\
&= \frac{\frac{1}{\rho_{\mathrm{target}}} - \frac{1}{\rho'}}
{\frac{1}{\rho''} - \frac{1}{\rho'}}.
\end{aligned}
\end{equation}
Именно это соотношение
связывает заданную среднюю плотность смеси
с массовой долей пара.

После вычисления $x$
аддитивные по массе свойства
определяются по обычному правилу смешения:
\begin{equation}
y = (1-x)\,y' + x\,y''.
\end{equation}
Если вычисленное значение $x$
выходит за диапазон $[0; 1]$,
то состояние не принадлежит двухфазной области
и должно быть возвращено на однофазную маршрутизацию.

Вблизи критической точки
разность $v''-v'$ стремится к нулю.
Поэтому формула для $x$
становится плохо обусловленной:
небольшая ошибка в плотности
даёт заметную ошибку в степени сухости.
В реализации это компенсируется
контролем диапазонов,
проверкой близости к критической точке
и отдельной обработкой пограничных случаев.


\section{Ограничения применимости и обработка особых случаев}

Корректность расчётов определяется
областью применимости IF97 и устойчивостью численных процедур.
Поэтому в ядре выполняются
валидация входных параметров и явная обработка ошибок.

К типовым ограничениям относятся:
\begin{itemize}
\item допустимые диапазоны давления и температуры в соответствии со стандартом IAPWS-IF97~\cite{iapws_if97_2012};
\item требование $\rho > 0$ для маршрутов, использующих плотность;
\item требование $x \in [0; 1]$ для расчётов в двухфазной области;
\item ограничение области метастабильного пара в расширении региона~2 по давлению.
\end{itemize}

К особым случаям относятся вычисления на фазовых границах, в околокритической области и при переходах между регионами.
В этих ситуациях важны:
\begin{itemize}
\item контроль смены региона при итерационном процессе;
\item обработка случаев плохой обусловленности (например, вблизи критической точки);
\item ограничение числа итераций и явная сигнализация о несходимости;
\item единый формат ошибок валидации и вычислений, пригодный для интерфейсного и сервисного использования.
\end{itemize}

Перечисленные меры обеспечивают
устойчивую работу модели в инженерных сценариях
и воспроизводимость результатов в \texttt{if97\_core},
FLTK,
в оболочке на базе Yew, WebAssembly и Tauri
и в gRPC-сервисе.

